\begin{textsample}{POS Dim 1 – rj – Score 31.00 – rj/rj\_ef\_linguagens\_lp\_1.txt}  \label{ex:f1_pos_252}
Língua Portuguesa O ensino de Língua Portuguesa precisa dar continuidade às práticas já iniciadas e consolidadas nos anos anteriores em seus diversos usos e contextos sem, \textbf{contudo}, \textbf{que} o professor permita \textbf{que} as variantes linguísticas de cada localidade sejam tratadas como inferiores.
Desde pequenos os estudantes precisam conhecer, analisar e discutir a diversidade linguística da localidade em \textbf{que} \textbf{vivem} e compreender a riqueza dessas variantes para os grupos sociais aos quais pertencem, cabendo ao professor de Língua Portuguesa esse papel tão importante.
Mesmo sabendo \textbf{que} o estudante deve ser alfabetizado em sua língua oficial, a autonomia pedagógica do professor permite \textbf{que} discuta com os mesmos sobre aspectos relevantes da cultura local, regional, nacional e mundial, desde \textbf{que} respeitada à faixa etária e a \textbf{maturidade} de cada grupo.
A presença de \textbf{um} trabalho pautado na diversidade textual requer \textbf{um} lugar de destaque no cotidiano escolar, pois através dele os estudantes iniciam o processo de autonomia, pois a leitura e a escrita são, incontestavelmente, o caminho para \textbf{que} estejam inseridos em nossa sociedade.
Durante o trabalho pedagógico com a Língua Portuguesa o enfoque do professor precisa ser o texto em seus variados contextos, o \textbf{que} também significa inserir as variantes locais em seu dia a dia escolar, de maneira \textbf{que} o texto, oral, escrito ou multimodal, passe a ser o centro das atividades a serem desenvolvidas no espaço escolar, permitindo \textbf{que} haja \textbf{uma} troca permanente de ideias, questionamentos e \textbf{posicionamentos}, respeitando assim as individualidades e necessidades dos envolvidos no processo.
A proposta pedagógica de cada Rede articulada com as necessidades dos estudantes pode resultar em \textbf{um} trabalho cujo foco será os variados gêneros textuais e diferentes aspectos gramaticais da Língua Portuguesa para \textbf{que} os estudantes percebam \textbf{que} a Língua vai muito além do \textbf{que} \textbf{um} \textbf{simples} sistema de códigos e regras gramaticais.
No ambiente de multiletramentos em \textbf{que} \textbf{vivemos} torna -se indispensável \textbf{que} o processo ensino-aprendizagem abra espaços para a reflexão e a criticidade, para \textbf{que} os estudantes se tornem cada vez mais autônomos em suas produções ( faladas e escritas ), utilizando recursos linguísticos adequados.
Nesse sentido, através de \textbf{um} trabalho interdisciplinar no qual o professor pode atuar como mediador, a produção de conhecimento pelos estudantes se dará através dos usos das variadas linguagens presentes nos diferentes tipos textuais utilizados em sala de aula, associadas às demandas de leitura e escrita advindas da prática social \textbf{que} \textbf{permeiam} o ambiente em \textbf{que} \textbf{vivem}, dentro e fora dela.
A escola precisa ser \textbf{um} local de interlocução, ou seja, precisa abrir espaços para o uso de diferentes linguagens, interações e discursos, ouvindo o \textbf{que} professores e estudantes têm a \textbf{dizer} a partir de suas experiências e vivências, já \textbf{que} ela é o lugar de \textbf{excelência} para o aprendizado.
Sendo a língua viva, entendida como algo inacabado, em constante transformação e \textbf{que} acontece no discurso, é de \textbf{suma} importância levar em consideração seus diferentes usos, em diferentes contextos e situações.
Para compreender e utilizar a linguagem corretamente os estudantes precisam pensar sobre ela.
Soares ( 2003 ) afirma \textbf{que} “ não se escreve como se fala, mesmo quando se fala em situações formais ; não se fala como se escreve, mesmo quando se escreve em contextos informais ”.
Considerando ainda \textbf{que} a ação pedagógica do professor está \textbf{permeada} por várias concepções e \textbf{que} estas concepções norteiam toda sua prática e autonomia no fazer pedagógico, é necessário \textbf{que} se tenha clareza sobre elas para \textbf{que} haja \textbf{uma} melhor organização do trabalho.
Segundo Travaglia ( 2001, p. 10-11 ) : > [... ] o professor deve evitar a adesão superficial a modismos linguísticos ou da \textbf{pedagogia} de língua materna, sem, pelo menos, \textbf{um} conhecimento substancial das \textbf{teorias} linguísticas em \textbf{que} se \textbf{embasam} e dos pressupostos de todos os tipos ( linguísticos, pedagógicos, psicológicos, políticos, etc. ) \textbf{que} dão forma a \textbf{teorias} e \textbf{métodos}.
A ansiedade de inovar ou parecer \textbf{moderno} nos leva muitas vezes a maquilar \textbf{teorias} e \textbf{métodos} antigos com aspectos superficiais de novas \textbf{teorias} e \textbf{métodos}, gerando não bons instrumentos de trabalho, mas verdadeiras degenerações \textbf{que} mais perturbam do \textbf{que} ajudam, por não se saber exatamente o \textbf{que} se está fazendo.
Daí \textbf{um} pressuposto óbvio de toda metodologia, mas no qual devemos insistir : não há bom ensino sem o conhecimento profundo do objeto de ensino ( no nosso caso, da Língua Portuguesa ) e dos elementos \textbf{que} dão forma ao \textbf{que} realizamos em sala de aula em função de muitas opções \textbf{que} fazemos ou \textbf{que} não fazemos. (... ) É \textbf{preciso}, pois, estar consciente das opções \textbf{que} fazemos (... ), ao estruturar e realizar o ensino de Português para falantes dessa língua, em \textbf{face} dos objetivos \textbf{que} se julgam pertinentes ( estes já são \textbf{uma} opção ) para se dar aulas de \textbf{uma} língua a seus falantes nativos.
Sendo assim, toda ação docente precisa estar pautada em objetivos claros \textbf{que} \textbf{explicitem} tanto o papel do professor como mediador, quanto o \textbf{que} esperam \textbf{que} os estudantes \textbf{consigam} atingir no decorrer do caminho e ao final de cada etapa do processo ensino-aprendizagem.
A pertinência da Língua Portuguesa à área de Linguagens é indiscutível, sendo ela por assim \textbf{dizer}, o elo interdisciplinar entre os demais componentes curriculares, de caráter total de abrangência : \textbf{remete} à expressão oral mais \textbf{simples}, ganha magia na alfabetização contextualizada, perpassa pelos gêneros textuais primários e secundários, evolui com o conhecimento gramatical, favorecendo a produção textual e o protagonismo de seus usuários.
É \textbf{preciso} \textbf{que} o professor tenha clareza de \textbf{que} aprender Língua Portuguesa vai muito além de ensinar aspectos gramaticais, itens lexicais, coesão e até mesmo apenas aspectos linguísticos.
A estudante \textbf{precisa} de propostas pedagógicas \textbf{que} os alcancem na maneira como aprendem, conhecimentos sobre e com a língua, precisam aprender sobre as diversas linguagens e suas variações, respeitando -as.
Absolutamente integrada a todas as outras áreas, a Língua Portuguesa assume de forma mais abrangente do \textbf{que} nunca o conceito de \textbf{interdisciplinaridade} para a formação integral do estudante.
Ser letrado é ter a capacidade de escutar, falar, ler e escrever sobre todos os assuntos presentes em qualquer componente curricular e nas questões pessoais ou coletivas, sempre atendendo às adaptações necessárias de comunicação, utilizando inclusive tecnologias digitais da informação e comunicação ( TDIC ), letramento digital e midiático.
Participar do mundo letrado é fundamental quando a educação é prevista para todos e dessa forma o estudante precisa saber escutar, falar, ler e escrever, compreendendo, interpretando, produzindo e observando os aspectos culturais de sua língua materna, suas variantes, seus conhecimentos linguísticos, gramaticais, os gêneros textuais e suas estruturas, mas sempre à luz da linguagem.
Desta maneira, os estudantes podem ampliar sua participação na vida escolar e nos grupos sociais aos quais pertencem e/ou convivem, possibilitando \textbf{que} participem de maneira cada vez mais autônoma, crítica e reflexiva para \textbf{que} possam contribuir de maneira significativa para possíveis mudanças em suas práticas sociais cotidianas.
É de \textbf{suma} importância \textbf{que} os estudantes também aprendam a lidar com os conteúdos \textbf{que} aparecem nas webs, pois ter tido contato e saber utilizar de maneira ética e consciente, é muito diferente.
A escola precisa acolher essas informações, apresentando, conversando e discutindo com o grupo sobre conteúdos \textbf{que} circulam nas redes sociais, de acordo com cada faixa etária.
Os estudantes precisam saber pesquisar fontes seguras, diferenciar o \textbf{que} é real do \textbf{que} é inventado, o \textbf{que} pode se tornar público ou não, o \textbf{que} o excesso de \textbf{exposição} e os discursos de ódio podem acarretar na vida das pessoas, políticas de uso e outras questões \textbf{que} aparecerão no decorrer de sua vida escolar, cabendo ao professor, imbuído de conhecimentos e práticas pedagógicas adequadas, mediar essa construção de conhecimentos no espaço escolar.
Cabe à escola aprender a lidar com a diversidade e as diferenças abrindo espaços para \textbf{que} os estudantes possam se colocar, debater e \textbf{argumentar} sobre esses assuntos, ajudando -os a lidar com as transformações e as novas demandas cada vez mais presente nos espaços sociais.
É fundamental \textbf{que} os estudantes discutam, reflitam, aprendam, entendam e compreendam sobre a responsabilidade pelo uso incorreto das TDIC, por isso discutir e definir parâmetros para seu uso é primordial, é de \textbf{suma} importância saber \textbf{que} existem regras e limites a serem respeitados nesse universo amplo de comunicação e expressão.
Os novos e multiletramentos \textbf{que} vieram surgindo com a sociedade e com o mundo \textbf{globalizado} no qual \textbf{vivemos} \textbf{hoje}, não permitem \textbf{que} tenhamos as mesmas práticas e respostas, pois embora as questões sejam praticamente as mesmas, precisamos buscar novas soluções em nossas salas de aula fazendo com \textbf{que} os estudantes também aprendam a buscá -las.
Fazer com \textbf{que} enfrentem os desafios do dia a dia de maneira criativa, autônoma, reflexiva, mas, sobretudo ética, também é \textbf{um} dos desafios da educação nos dias \textbf{atuais}.
É \textbf{preciso} se pensar em como essas novas tecnologias da informação, esses hipertextos e essas hipermídias podem contribuir na maneira com \textbf{que} professores e estudantes se relacionam nas salas de aula, em como podem contribuir positivamente no processo ensino-aprendizagem, já \textbf{que} muitas profissões ainda surgirão no decorrer do tempo e precisarão de diferentes habilidades para exercê -las, por isso o trabalho com a cultura digital é de extrema importância nas atividades pedagógicas dos dias \textbf{atuais}.
A diversidade precisa ser respeitada e entendida para \textbf{que} todos possam usufruir o direito à igualdade e, para \textbf{que} isso se torne realidade, são necessários planos de aula \textbf{que} tragam a diversidade à pauta para \textbf{que} estudantes e professores se conscientizem de \textbf{que} o nosso país é composto por misturas \textbf{que} merecem ser respeitadas e valorizadas, pois todos têm o mesmo valor diante da sociedade, independente das \textbf{singularidades}, inclusive a variedade linguística, de cada \textbf{um}.
No momento da construção dos planos de aulas o protagonismo dos professores, a \textbf{interdisciplinaridade} e o estudo \textbf{teórico} das metodologias linguísticas deverão vir à tona para \textbf{que} se chegue a \textbf{um} dos objetivos \textbf{centrais} das aulas de Língua Materna \textbf{que} é letrar os estudantes em sua língua fazendo com \textbf{que} os mesmos possam desenvolver a linguagem apropriada ao seu uso em diferentes contextos sociais.

% matched lemmas: argumentar, atual, central, conseguir, contudo, dizer, embasar, excelência, explicitar, exposição, face, globalizado, hoje, interdisciplinaridade, maturidade, moderno, método, pedagogia, permear, posicionamento, preciso, que, remeter, simples, singularidade, sumo, teoria, teórico, um, vivar, viver
\end{textsample}
